\section{Limites de la démarche}\label{sec:limits}

Cette exécution montre que le pipeline fonctionne, mais aussi par où il pèche. Nous passons ses limites en revue dans l'ordre des étapes de la figure~\ref{fig:pipeline}.

\paragraph{Les arches sont classiques.} Les cinq arches de \lp{} dans lean-pool portent sur des topos de \emph{préfaisceaux}: ce sont des résultats classiques de Caramello, Johnstone, Rogers et d'autres, qui ramènent un invariant à une condition sur une petite catégorie. La section~\ref{sec:analysis} emploie deux arches de plus, les caractérisations classiques des topos de faisceaux booléens et de De Morgan d'un espace; elles y sont énoncées, mais non formalisées. Des conditions sur une petite catégorie sont faciles à énoncer et à vérifier, et la plupart de nos résultats ont une preuve directe courte. Font exception le carquois cyclique (une page), $(\mathbb Z/n,\cdot)$ (la classification de Pirashvili et une factorisation de type Ramsey) et l'énoncé sur les compacts séparés (Stone et Nakano), dont les preuves utilisent des théorèmes cités. La profondeur d'un théorème engendré est bornée par celle de son arche et de son tablier \cite[transparent 8]{caramello-mines}; pour obtenir des résultats plus profonds, il faut des arches sur de vrais sites (par exemple les caractérisations d'invariants géométriques sur les sites dues à Caramello \cite{caramello-tst}) et des tabliers qui relient des théories équivalentes de façon non triviale. Notre registre ne contient ni arche sur de vrais sites ni tablier du second genre; la note associée \cite{bridges-notes} en traite un, le topos des pseudo-corps tiré des notes de Grothendieck.

\paragraph{Les tabliers sont choisis à la main.} Les résultats nouveaux viennent de tabliers qui n'avaient pas encore été associés à ces arches: les représentations ensemblistes de carquois et les monoïdes de transformations classiques. Les choisir était l'étape créative, et rien dans le pipeline ne l'automatise. Un registre de la forme \q{la structure $X$ est un préfaisceau sur la forme $C$} est facile à remplir pour les topos de préfaisceaux; pour les topos classifiants de théories, le remplir systématiquement demande de démontrer des équivalences de Morita, ce qui relève de la recherche.

\paragraph{Spécialiser peut être difficile.} Calculer la condition sur la forme était facile pour Ore sur $\B M$ ou sur une catégorie libre, mais a déjà demandé une page pour les topologies sur le carquois cyclique (théorème~\ref{th:cycle}); les topologies sur $\B(\mathbb N^2)$, ou la condition d'Ore sur une catégorie présentée par générateurs et relations, sont de vrais problèmes. Le générateur déplace la difficulté sans la supprimer.

\paragraph{Déplier et réécrire demandent un humain.} La sortie brute est \q{$\mathrm{Paths}(R^{\op})$ vérifie \texttt{RightOreCondition}}; l'énoncé que veut un mathématicien est \q{chaque sommet a au plus une flèche sortante}. C'est cette étape que l'exposé confie aux systèmes de réécriture et aux modèles de langage \cite[transparent 17]{caramello-mines}, et c'est elle qui nous a pris la plus grande partie de notre temps.

\paragraph{Une exécution automatique des premières étapes.} Pour éprouver la tactique que propose l'exposé sur les deux bases de données \cite[transparent 17]{caramello-mines}, nous avons mis les arches booléenne et de De Morgan, et leurs lectures sur les tabliers \texttt{SingleObj M} et \texttt{Paths Q}, dans un jeu de règles Aesop \cite{aesop}, comme règles de normalisation (\lean{lean/bridges/Bridges/Aesop/}). Devant les énoncés déjà démontrés à la main, \texttt{aesop} trouve la preuve de ceux-ci, et de ceux-ci seulement:
\begin{itemize}
\item \emph{joindre et spécialiser} (étapes~1 et~2): les quatre lectures des deux arches sur les deux tabliers, en un appel chacune;
\item \emph{conclure par un fait connu}: le topos booléen d'un groupe, à condition de donner \texttt{Group.isUnit} comme règle;
\item \emph{un transfert de Morita} (un groupoïde reste booléen le long d'une équivalence de topos de préfaisceaux), à condition d'énoncer l'arche à l'envers dans un lemme à part, de déplier l'équivalence et de donner l'invariance de la propriété comme règle vers l'avant.
\end{itemize}
Il ne trouve pas les corollaires dont la condition sur le tablier demande un témoin: De Morgan pour les monoïdes commutatifs ($c=b$, $d=a$) et pour les monoïdes avec zéro ($c=d=0$), ni l'échec de De Morgan pour le monoïde libre sur deux lettres. Chaque échec est consigné dans le fichier sous \texttt{fail\_if\_success}, de sorte que la compilation casse si la frontière bouge. Deux limites de l'outil apparaissent en chemin: une règle \texttt{simp} ne s'emploie pas de droite à gauche, il faut donc écrire l'arche renversée; et une recherche qui échoue laisse le but ouvert, il faut donc lancer les échecs en recherche terminale. L'automatisation couvre ainsi les étapes qui étaient déjà mécaniques; les témoins, les transferts et le choix des tabliers restent à l'humain. Elle ne propose aucun énoncé: elle démontre ceux qu'on lui donne.

\paragraph{Lean vérifie le corollaire, pas sa lecture.} Une preuve Lean que $\mathrm{Paths}(R^{\op})$ vérifie Ore si et seulement si une condition est satisfaite est une garantie sur cet énoncé, dans les conventions de la bibliothèque (côté de l'action, sens des flèches, \texttt{SingleObj} composant selon $f\gg g = g\cdot f$). Que la phrase informelle de l'article dise la même chose, Lean ne le vérifie pas. Nous avons rencontré deux conventions qui changent le côté d'une action de monoïde, et ce côté compte: les $T_n$-ensembles à droite ($ab = a\circ b$) ne forment pas un topos de De Morgan pour $n\ge 2$, puisque les applications constantes $c_i$ vérifient $c_iT_n = \{c_i\}$, alors que les $T_n$-ensembles à gauche en forment un.

\paragraph{Le filtre de nouveauté est faible.} \q{Non trouvé} ne veut pas dire \q{nouveau}. Le jeu d'appariements de l'étape~6 n'atteint que les énoncés qui ont un jumeau Lean. Notre recherche (annexe~\ref{app:novelty}) a couvert les sources évidentes, mais les énoncés élémentaires relèvent souvent du folklore: on les rencontre en passant ou en exercice, et ils ne sont pas indexés. Le journal d'étalonnage compte 14 énoncés, contre 10 lignes dans le tableau des résultats, et trois énoncés que nous avions crus nouveaux se sont révélés publiés (le critère pour les monoïdes finis, les sous-topos des $\mathbb N$-ensembles, les points essentiels comme classes d'idempotents). Nous n'avons pas consigné le nombre de couples (tablier, arche) essayés (\S\ref{sec:calibration}): ces chiffres ne donnent donc aucune proportion, mais ils montrent que le filtre de nouveauté a changé ce que nous rapportons.

\paragraph{La plupart des sorties sont sans intérêt.} Associer chaque tablier à chaque arche produit beaucoup d'énoncés vrais, nouveaux au sens étroit, et sans aucun intérêt (le tiers exclu faible pour les préfaisceaux sur une catégorie particulière à cinq objets). L'avertissement de l'exposé sur les résultats \q{bizarres} s'applique ici. Pour être utile, un générateur doit classer ses sorties; cela relève du goût mathématique, et ce n'est pas formalisé.

\paragraph{À quoi sert la démarche.} Dans ces limites, le pipeline sert à quelque chose: il pose des questions que personne n'avait posées (le tiers exclu faible est-il vrai pour les automates? combien de sous-topos un cycle d'applications a-t-il?), il y répond avec une garantie (l'arche est démontrée formellement) et il indique une preuve directe. Comme outil d'enseignement, ou pour explorer la frontière entre deux domaines, c'est déjà appréciable.

\subsection{Leçons pour un générateur automatique}

\begin{enumerate}
\item \emph{Choisir les tabliers est l'étape créative.} Les résultats nouveaux viennent de tabliers que nous n'avons pas trouvés associés à ces arches: carquois et monoïdes de transformations. Une base de données de tabliers utile est une base d'énoncés de la forme \q{cette structure familière est un préfaisceau sur telle forme}.
\item \emph{L'étape de dépliage est mécanique.} Le lemme~\ref{lem:neg} est la seule règle nécessaire pour traduire \q{De Morgan} en un énoncé sur des états et des pièges. Un générateur a besoin d'une telle règle par invariant.
\item \emph{C'est la réécriture qui demande le travail.} La sortie brute pour les automates est \q{$\B(\Sigma^*)$ vérifie \texttt{RightOreCondition}}; l'énoncé lisible porte sur les flèches sortantes d'un sommet (théorème~\ref{th:quivers}). C'est l'étape que l'exposé confie aux systèmes de réécriture et aux modèles de langage \cite[transparent 17]{caramello-mines}.
\item \emph{Le filtrage par la littérature est nécessaire.} Ce n'est qu'à ce stade que nous avons trouvé publiés trois de nos candidats (\S\ref{sec:calibration}). Faute d'avoir consigné le nombre de couples essayés, nous ne pouvons donner aucune proportion. Sans filtre de nouveauté, un générateur présente des résultats connus comme nouveaux.
\item \emph{La vérification.} Chaque résultat présenté ici a une preuve directe: courte pour la plupart, d'une page pour le carquois cyclique, plus longue pour $(\mathbb Z/n,\cdot)$ et pour l'énoncé sur les compacts séparés, qui utilisent des théorèmes cités (la classification de Pirashvili, avec une factorisation de type Ramsey; Stone et Nakano). Pour des tabliers plus profonds, une preuve directe peut ne pas exister, et la preuve formelle de l'arche en Lean sera alors la seule garantie. C'est pourquoi on formalise d'abord les arches.
\end{enumerate}
